\documentclass[a4paper,12pt]{article}
\usepackage[a4paper, margin=1.8cm]{geometry}
\usepackage{tikz}
\usepackage{amsmath}
\usepackage{amssymb}
\usepackage{array}
\usepackage{xcolor}
\usepackage{mdframed}
\usepackage{enumitem}
\usepackage[T1]{fontenc}
\usepackage{textcomp}
\usetikzlibrary{calc}

% ---- Colours ----
\definecolor{slice1}{RGB}{70,130,180}
\definecolor{slice2}{RGB}{240,140,40}
\definecolor{slice3}{RGB}{70,170,70}
\definecolor{slice4}{RGB}{210,70,70}
\definecolor{slice5}{RGB}{150,90,190}
\definecolor{lightgrey}{RGB}{245,245,245}
\definecolor{boxblue}{RGB}{220,235,250}
\definecolor{boxgreen}{RGB}{220,245,220}

% -------------------------------------------------------
% \piechart{radius_number}{Label/Angle/Colour, ...}
%   radius_number is a plain number interpreted as cm.
%   Angles in degrees, clockwise from 12 o'clock; must sum to 360.
% -------------------------------------------------------
\newcommand{\piechart}[2]{%
  \begin{tikzpicture}
    \pgfmathsetmacro{\r}{#1}%
    \pgfmathsetmacro{\startA}{90}%
    \foreach \lbl/\ang/\col in {#2} {%
      \pgfmathsetmacro{\endA}{\startA - \ang}%
      \fill[\col!20] (0,0) -- (\startA:\r cm) arc(\startA:\endA:\r cm) -- cycle;%
      \draw[line width=2pt] (0,0) -- (\startA:\r cm);%
      \pgfmathsetmacro{\midA}{(\startA + \endA)/2}%
      \pgfmathsetmacro{\lR}{\r * 0.65}%
      \node[font=\small\bfseries, align=center] at (\midA:\lR cm) {\lbl};%
      \xdef\startA{\endA}%
    }%
    \draw[black, line width=1.8pt] (0,0) circle (\r cm);%
    \draw[black, line width=2pt] (0,0) -- (90:\r cm);%
    \fill[black] (0,0) circle (3pt);%
    \draw[black, line width=1.2pt] (0,\r cm) -- (0,{\r cm+0.25cm});%
  \end{tikzpicture}%
}

% Shortcut for a boxed written answer
\newcommand{\ans}[1]{\fbox{\textbf{#1}}}

\pagestyle{empty}

\begin{document}

% ===========================
%  TITLE
% ===========================
\begin{center}
  {\LARGE\bfseries Reading Pie Charts with a Protractor}\\[3pt]
  {\large Worked Solutions}
\end{center}
\vspace{4pt}

\begin{mdframed}[backgroundcolor=boxgreen, linecolor=slice3,
                 linewidth=1.5pt, roundcorner=6pt, innertopmargin=6pt]
\textbf{Method:}\quad
$\displaystyle\text{sector fraction}=\frac{\text{sector angle}}{360}$
\qquad
$\displaystyle\text{amount}=\text{fraction}\times\text{total}$

\smallskip
For every sector, use the exact angle shown on the chart, then simplify the fraction where possible.
\end{mdframed}

\newpage

% ===========================
%  PIE CHART 1
% ===========================
\begin{center}
  {\large\bfseries Pie Chart 1 \textendash{} Favourite School Subjects (30 students)}
\end{center}

\noindent
\begin{minipage}[c]{0.6\textwidth}
  \centering
  \piechart{5.8}{Maths/120/slice1,Science/90/slice2,English/72/slice3,Art/48/slice4,PE/30/slice5}
\end{minipage}%
\hfill
\begin{minipage}[c]{0.36\textwidth}
  \setlength{\tabcolsep}{4pt}
  \renewcommand{\arraystretch}{1.6}
  \begin{tabular}{|l|c|c|c|}
  \hline
  \textbf{Subject} & \textbf{Angle} & \textbf{Fraction} & \textbf{No.} \\
  \hline
  Maths    & \(120^\circ\) & \(\frac{120}{360}=\frac13\) & \(10\) \\
  \hline
  Science  & \(90^\circ\) & \(\frac14\) & \(7.5^*\) \\
  \hline
  English  & \(72^\circ\) & \(\frac15\) & \(6\) \\
  \hline
  Art      & \(48^\circ\) & \(\frac{2}{15}\) & \(4\) \\
  \hline
  PE       & \(30^\circ\) & \(\frac{1}{12}\) & \(2.5^*\) \\
  \hline
  \textbf{Total} & \textbf{360\textdegree} & \textbf{1} & \textbf{30} \\
  \hline
  \end{tabular}
\end{minipage}

\vspace{4pt}
\noindent\textit{Note:} \(^*\)The exact angles shown give fractional numbers of students for Science and PE. This is a reminder that a pie chart must be designed carefully if the data are counts of people.

\vspace{4pt}
\noindent\textbf{Worked Solutions -- Pie Chart 1:}
\begin{enumerate}[leftmargin=1.6em, itemsep=4pt]
  \item Measuring the sectors gives the angles in the table above:
        \[
        120^\circ,\;90^\circ,\;72^\circ,\;48^\circ,\;30^\circ.
        \]
        Their sum is
        \[
        120+90+72+48+30=360^\circ.
        \]

  \item The angles add up to \(\boxed{360^\circ}\).

  \item Maths fraction:
        \[
        \frac{120}{360}=\boxed{\frac13}.
        \]

  \item Science students:
        \[
        \frac{90}{360}\times 30
        =\frac14\times 30
        =\boxed{7.5\text{ students}}.
        \]

  \item A fraction of \(\frac15\) of the circle is
        \[
        \frac15\times 360=72^\circ.
        \]
        The \(72^\circ\) sector is English, so the subject is \(\ans{English}\).
\end{enumerate}

\newpage

% ===========================
%  PIE CHART 2
% ===========================
\begin{center}
  {\large\bfseries Pie Chart 2 \textendash{} How Students Travel to School (60 students)}
\end{center}

\noindent
\begin{minipage}[c]{0.6\textwidth}
  \centering
  \piechart{5.8}{Walk/90/slice3,Bus/120/slice1,Car/60/slice2,Cycle/45/slice4,Other/45/slice5}
\end{minipage}%
\hfill
\begin{minipage}[c]{0.36\textwidth}
  \setlength{\tabcolsep}{4pt}
  \renewcommand{\arraystretch}{1.6}
  \begin{tabular}{|l|c|c|c|}
  \hline
  \textbf{Method} & \textbf{Angle} & \textbf{Fraction} & \textbf{No.} \\
  \hline
  Walk   & \(90^\circ\) & \(\frac14\) & \(15\) \\
  \hline
  Bus    & \(120^\circ\) & \(\frac13\) & \(20\) \\
  \hline
  Car    & \(60^\circ\) & \(\frac16\) & \(10\) \\
  \hline
  Cycle  & \(45^\circ\) & \(\frac18\) & \(7.5^*\) \\
  \hline
  Other  & \(45^\circ\) & \(\frac18\) & \(7.5^*\) \\
  \hline
  \textbf{Total} & \textbf{360\textdegree} & \textbf{1} & \textbf{60} \\
  \hline
  \end{tabular}
\end{minipage}

\vspace{4pt}
\noindent\textit{Note:} \(^*\)Cycle and Other each give exactly \(7.5\) students. If whole students are required, each is approximately \(8\), but the exact calculation using the sector angles is \(7.5\).

\vspace{4pt}
\noindent\textbf{Worked Solutions -- Pie Chart 2:}
\begin{enumerate}[leftmargin=1.6em, itemsep=4pt]
  \item The measured sector angles are:
        \[
        90^\circ,\;120^\circ,\;60^\circ,\;45^\circ,\;45^\circ,
        \]
        and their sum is
        \[
        90+120+60+45+45=360^\circ.
        \]

  \item Students travelling by Bus:
        \[
        \frac{120}{360}\times 60
        =\frac13\times 60
        =\boxed{20\text{ students}}.
        \]

  \item Fraction of students who walk:
        \[
        \frac{90}{360}=\boxed{\frac14}.
        \]

  \item The two sectors with the same angle are:
        \[
        45^\circ\text{ and }45^\circ,
        \]
        so the two methods are \(\ans{Cycle and Other}\).

  \item Bus minus Car:
        \[
        20-10=\boxed{10\text{ more students}}.
        \]
\end{enumerate}

\newpage

% ===========================
%  PIE CHART 3
% ===========================
\begin{center}
  {\large\bfseries Pie Chart 3 \textendash{} Ice Cream Flavours Sold in One Day (180 sold)}
\end{center}

\noindent
\begin{minipage}[c]{0.6\textwidth}
  \centering
  \piechart{5.8}{Vanilla/80/slice2,Choc/100/slice1,Straw/60/slice4,Mint/40/slice3,Other/80/slice5}
\end{minipage}%
\hfill
\begin{minipage}[c]{0.36\textwidth}
  \setlength{\tabcolsep}{4pt}
  \renewcommand{\arraystretch}{1.6}
  \begin{tabular}{|l|c|c|c|}
  \hline
  \textbf{Flavour} & \textbf{Angle} & \textbf{Fraction} & \textbf{No.} \\
  \hline
  Vanilla    & \(80^\circ\) & \(\frac{2}{9}\) & \(40\) \\
  \hline
  Chocolate  & \(100^\circ\) & \(\frac{5}{18}\) & \(50\) \\
  \hline
  Strawberry & \(60^\circ\) & \(\frac16\) & \(30\) \\
  \hline
  Mint       & \(40^\circ\) & \(\frac19\) & \(20\) \\
  \hline
  Other      & \(80^\circ\) & \(\frac{2}{9}\) & \(40\) \\
  \hline
  \textbf{Total} & \textbf{360\textdegree} & \textbf{1} & \textbf{180} \\
  \hline
  \end{tabular}
\end{minipage}

\vspace{4pt}
\noindent\textbf{Worked Solutions -- Pie Chart 3:}
\begin{enumerate}[leftmargin=1.6em, itemsep=4pt]
  \item The sector angles are:
        \[
        80^\circ,\;100^\circ,\;60^\circ,\;40^\circ,\;80^\circ,
        \]
        and
        \[
        80+100+60+40+80=360^\circ.
        \]

  \item Chocolate ice creams sold:
        \[
        \frac{100}{360}\times 180
        =\frac{5}{18}\times 180
        =\boxed{50\text{ ice creams}}.
        \]

  \item Mint ice creams sold:
        \[
        \frac{40}{360}\times 180
        =\frac19\times 180
        =\boxed{20\text{ ice creams}}.
        \]

  \item Vanilla and Other both have angle \(80^\circ\). Each represents:
        \[
        \frac{80}{360}\times 180
        =\frac{2}{9}\times 180
        =\boxed{40\text{ ice creams each}}.
        \]

  \item A flavour that is exactly \(\frac16\) of all sales has angle
        \[
        \frac16\times 360=60^\circ.
        \]
        The \(60^\circ\) sector is Strawberry.
        \[
        \frac{60}{360}\times 180=30,
        \qquad
        \frac16\times 180=30.
        \]
        So the flavour is \(\ans{Strawberry}\).
\end{enumerate}

\newpage

% ===========================
%  PIE CHART 4
% ===========================
\begin{center}
  {\large\bfseries Pie Chart 4 \textendash{} One Evening's Activities (4 hours = 240 minutes)}
\end{center}

\noindent
\begin{minipage}[c]{0.6\textwidth}
  \centering
  \piechart{5.8}{HW/90/slice4,TV/120/slice2,Gaming/60/slice1,Reading/30/slice3,Other/60/slice5}
\end{minipage}%
\hfill
\begin{minipage}[c]{0.36\textwidth}
  \setlength{\tabcolsep}{4pt}
  \renewcommand{\arraystretch}{1.6}
  \begin{tabular}{|l|c|c|c|}
  \hline
  \textbf{Activity} & \textbf{Angle} & \textbf{Fraction} & \textbf{Mins} \\
  \hline
  Homework & \(90^\circ\) & \(\frac14\) & \(60\) \\
  \hline
  TV       & \(120^\circ\) & \(\frac13\) & \(80\) \\
  \hline
  Gaming   & \(60^\circ\) & \(\frac16\) & \(40\) \\
  \hline
  Reading  & \(30^\circ\) & \(\frac{1}{12}\) & \(20\) \\
  \hline
  Other    & \(60^\circ\) & \(\frac16\) & \(40\) \\
  \hline
  \textbf{Total} & \textbf{360\textdegree} & \textbf{1} & \textbf{240} \\
  \hline
  \end{tabular}
\end{minipage}

\vspace{4pt}
\noindent\textbf{Worked Solutions -- Pie Chart 4:}
\begin{enumerate}[leftmargin=1.6em, itemsep=4pt]
  \item The measured angles are:
        \[
        90^\circ,\;120^\circ,\;60^\circ,\;30^\circ,\;60^\circ,
        \]
        and
        \[
        90+120+60+30+60=360^\circ.
        \]
        Since \(4\) hours \(=4\times 60=240\) minutes, each fraction is multiplied by \(240\).

  \item Time spent watching TV:
        \[
        \frac{120}{360}\times 240
        =\frac13\times 240
        =\boxed{80\text{ mins}}.
        \]

  \item Time spent on Homework:
        \[
        \frac{90}{360}\times 240
        =\frac14\times 240
        =\boxed{60\text{ mins}}.
        \]

  \item Fraction of the evening spent Gaming:
        \[
        \frac{60}{360}=\boxed{\frac16}.
        \]

  \item \textbf{Challenge:} Homework is \(\frac14\) of the evening.
        For another student with \(5\) hours \(=300\) minutes:
        \[
        \frac14\times 300
        =\boxed{75\text{ mins}}.
        \]
\end{enumerate}

\end{document}